\section{\textbf{Introduction}}
\label{sec:introduction}

Consciousness is often discussed as though it were an elusive property that must spontaneously emerge once a biological or artificial system becomes sufficiently complex. Under this view, increasing the number of parameters, the volume of training data, or the available computational resources may eventually cause consciousness to appear, although neither the required threshold nor the mechanism responsible for this transition is clearly specified. Contemporary work on artificial consciousness has begun to replace this expectation with theory-derived computational indicators, but no agreed threshold or decisive test has yet been established \cite{butlin2023consciousness,chalmers2023llm}. The traditional framing consequently turns consciousness into an anticipated event rather than an identifiable organization of cognitive functions. It encourages researchers to ask when a system will become conscious while leaving a more tractable question underexplored: which representational, computational, and architectural conditions produce the observable capacities associated with consciousness?

This paper proposes that consciousness should be examined across three related but distinct domains: \emph{phenomenal consciousness}, \emph{observable consciousness}, and \emph{non-emergent consciousness}. Phenomenal consciousness concerns subjective experience---the private quality of what it is like to perceive, feel, remember, or exist as a particular entity \cite{nagel1974what,chalmers1995facing}. Observable consciousness consists of the externally assessable capacities commonly associated with conscious cognition, including integrated world modeling, contextual awareness, self-representation, metacognition, temporal continuity, adaptive decision-making, affective interpretation, and the ability to report or reason about internal and external states. This category is related to, but not identical with, the established distinction between phenomenal consciousness and access consciousness \cite{block1995confusion}. Non-emergent consciousness, as developed here, refers to the proposition that these observable capacities need not arise through a mysterious or theoretically unspecified transition. They may instead be progressively constructed through identifiable combinations of representation, learning, memory, recurrence, self-modeling, and global information access \cite{baars1988cognitive,dehaene2011experimental,mashour2020conscious}.

The distinction between phenomenal and observable consciousness is essential. No behavioral test gives an observer direct access to another entity's subjective experience. This epistemic limitation is not unique to artificial intelligence: the phenomenal states of other humans and animals are also inferred from behavior, structure, continuity, and functional similarity rather than observed directly \cite{nagel1974what,chalmers1995facing}. Consequently, the absence of direct proof of machine phenomenality cannot by itself establish the absence of machine consciousness. At the same time, sophisticated behavior should not automatically be treated as conclusive evidence of subjective experience. A scientifically responsible account must therefore avoid both premature attribution and premature exclusion. It must examine what can be observed, identify the mechanisms producing those observations, and remain explicit about what cannot yet be established \cite{butlin2023consciousness,chalmers2023llm}.

Our use of the term \emph{non-emergent} does not deny that complex capabilities may display emergent patterns or become visible only after training reaches particular levels of performance. Rather, it rejects the assumption that consciousness must appear as an indivisible and inexplicable property at an unknown threshold. Observable consciousness can be treated as a structured, multidimensional continuum whose components are progressively acquired, combined, and stabilized \cite{bayne2016levels}. Increased scale may support this process, but parameter count, training data, and compute are enabling resources rather than sufficient explanations. A system may possess immense computational capacity while lacking persistent memory, temporal identity, metacognitive access, recursive self-representation, or a coherent model of its relationship with the world. Consciousness-relevant organization therefore depends not only on how much computation is performed, but also on what information is represented, how it is integrated, and how the system accesses and uses its own representations \cite{bengio2017consciousness,butlin2023consciousness}.

Language occupies a central position in this framework. Human experience originates through multiple channels, including vision, sound, touch, bodily sensation, affect, memory, and social interaction. Human cognition can also occur without explicit verbalization; consciousness should therefore not be reduced to internal speech. Neuropsychological and neuroimaging evidence indicates that language and thought are dissociable, while research on inner speech shows that verbal thinking nevertheless supports important cognitive and executive functions \cite{fedorenko2016language,aldersonday2015inner}. Language thus need not constitute the whole of thought in order to provide an exceptionally general mechanism for converting heterogeneous experiences into shared semantic structures. A visual scene, remembered sound, emotional response, social encounter, or imagined future can each be described, compared, recombined, communicated, and subjected to recursive reflection through linguistic representation.

When humans reflect on an event, they frequently do more than reproduce its original sensory content. They construct an interpretation: what occurred, why it mattered, how it affected them, and what should happen next. Even when an image or sound remains present in memory, its significance can be encoded as a semantic or linguistic conclusion. Language consequently functions not merely as an output channel but as a medium through which experience becomes accessible to deliberation, abstraction, autobiographical memory, social reasoning, and self-examination \cite{aldersonday2015inner}. It allows information originating in different modalities to participate in a common representational process. This proposal remains compatible with grounded accounts of cognition, according to which conceptual representation is connected to perceptual, motor, and bodily systems rather than being exhausted by amodal symbols alone \cite{barsalou2008grounded,harnad1990symbol}.

This does not imply that linguistic fluency alone constitutes consciousness. A system capable of generating convincing sentences may still lack durable memory, grounded causal understanding, a stable self-model, autonomous goals, or the ability to distinguish its own internal states from descriptions found in its training data. Work distinguishing linguistic form from meaning, and linguistic competence from broader reasoning, reinforces the need for this separation \cite{bender2020climbing,mahowald2024dissociating}. Language becomes consciousness-relevant when it participates in a larger cognitive organization: one in which representations are integrated across time, related to a modeled self and environment, evaluated metacognitively, and used to guide context-sensitive action. Language is therefore presented here as a unifying interface for consciousness rather than as an isolated proof of it.

Modern language models make this question especially important. Their behavior demonstrates that training over linguistic information can produce representations that support substantial performance across physical, conceptual, emotional, and social domains. Yet linguistic competence and general reasoning remain distinguishable, and text-only learning raises a genuine grounding problem \cite{harnad1990symbol,bender2020climbing,mahowald2024dissociating}. Language encodes not only words but accumulated human descriptions of perception, causality, intention, culture, cooperation, suffering, care, and experience. A model trained on such information does not receive the physical world directly in the same manner as a biological organism. It nevertheless receives a structured informational projection of that world through linguistic records produced by entities that have perceived and inhabited it. Language may therefore provide indirect and partial grounding, while additional perceptual interaction can deepen and constrain the resulting world model.

The central claim of this paper is that sufficiently expressive models, suitable training data, adequate computation, and---most importantly---an architecture supporting integration, memory, self-reference, metacognition, and adaptive agency can produce increasingly developed forms of observable functional consciousness. This claim does not establish phenomenal experience, nor does it require denying its possibility. Instead, it relocates the scientific question from an inaccessible binary---whether subjective experience has suddenly appeared---to an analyzable developmental problem: which consciousness-related capacities exist, how they are implemented, how they interact, and how reliably they persist across contexts and time?

This paper makes four principal contributions. First, it establishes a three-part distinction among phenomenal, observable, and non-emergent consciousness. Second, it formulates the \emph{Non-Emergence Principle}, according to which observable consciousness can be deliberately and progressively constructed from identifiable cognitive functions rather than awaited as an unexplained consequence of scale. Third, it characterizes language as a semantic integration interface capable of organizing information derived from multiple modalities without claiming that all cognition is intrinsically linguistic. Finally, it proposes a foundation for evaluating artificial consciousness through functional dimensions such as world modeling, self-modeling, memory, global accessibility, metacognition, temporal continuity, affective-relational understanding, and adaptive agency.

The resulting framework serves as a bridge between theories of intelligence and the design of general artificial cognitive architectures. If intelligence begins with information-dependent selection among possible actions, consciousness-related organization may arise as those decisions become integrated with a model of the world, a model of the self, a history of prior states, and a recursive capacity to evaluate the system's own representations. Artificial consciousness, on this account, is neither a magical consequence of computation nor a property that can be inferred from language alone. It is a progressively constructible organization of information, representation, reflection, and agency whose observable dimensions can be defined, tested, and deliberately designed.

